\section{Event Study: Permits}
\label{sec:empirical_results_permits}

\subsection{Redistricting Design}

The border discontinuity evidence in Section~\ref{sec:empirical_results_density} is strong, but one concern remains: opposite sides of ward borders could still differ in ways that border-segment fixed effects do not capture. I therefore exploit the 2015 ward redistricting as a natural experiment, which reassigned census blocks to new aldermen for reasons unrelated to housing market conditions. This lets me compare the same census block before and after its alderman changes, so block fixed effects absorb any time-invariant differences across locations.

Ward boundaries in Chicago are redrawn following each decennial census to equalize populations across wards. The 2015 redistricting (based on the 2010 Census) reassigned many census blocks to new aldermen while leaving nearby blocks under the same representation. This creates a difference-in-differences comparison between blocks that switched aldermen and nearby blocks that stayed in the same ward. The treatment is directional, not merely binary: some switched blocks moved to more stringent aldermen and some to more lenient ones. In the main specification, treatment enters continuously as the destination alderman's score minus the origin alderman's score, so positive values mean a move to a stricter alderman and negative values mean a move to a more lenient one. The identifying assumption is that, absent the change in aldermanic representation, treated and control blocks would have followed parallel trends in permit activity.

Several features of the design support this assumption. Redistricting is driven by population equalization and bargaining over demographic composition of the wards, not by housing market conditions on specific blocks. To tighten the counterfactual, I restrict the control group to blocks that remained in the \textit{origin} ward, so treated and control blocks shared the same pre-treatment regulatory environment. I also drop non-switching blocks in origin wards where the alderman changed between 2014 and 2015. Those would be contaminated controls: they did not redistrict into a new ward, but they still experienced an alderman change through election. Table~\ref{tab:redistricting_sample_support} reports the treatment and control group counts for the permit and home-sales redistricting designs. Figure~\ref{fig:identification_example} illustrates the design for one ward pair, showing how blocks are classified into treatment and control groups based on the change in alderman stringency. Figure~\ref{fig:treatment_control_maps} shows the citywide distribution of treated and control blocks for the 2015 redistricting.
Appendix~\ref{sec:appendix_event_study} reports treated-control balance for the permit event-study sample using block-level averages of exact parcel-based amenity distances and zoned FAR.

\input{../tasks/event_study_exact_balance/output/redistricting_sample_support_summary_250m.tex}

\begin{figure}[htbp]
    \centering
    \includegraphics[width=0.80\textwidth]{../tasks/event_study_sales_diagnostics/output/ward_pair_vertical_13-23_250m.pdf}
    \caption{Identification Example: Ward Pair 13--23 (2015 Redistricting)}
    \figurenotes{Top panel shows ward assignments before redistricting (2014 boundaries), middle panel shows assignments after redistricting (2015 boundaries), and bottom panel shows treatment status. Blocks that switch wards are treated, with treatment intensity given by the change in alderman stringency between the destination and origin wards. Blocks that remain in the origin ward serve as controls. Sample restricted to blocks within 250m of ward boundaries.}
    \label{fig:identification_example}
\end{figure}

\begin{figure}[htbp]
    \centering
    \includegraphics[width=0.62\textwidth]{../tasks/event_study_sales_diagnostics/output/treatment_control_map_2015_250m.pdf}
    \caption{Treatment and Control Blocks: 2015 Redistricting}
    \figurenotes{Geographic distribution of treated and control blocks for the 2015 redistricting. Red indicates blocks that moved to a more stringent alderman, blue indicates blocks that moved to a more lenient alderman, and gray indicates control blocks that did not switch wards. Sample restricted to blocks within 250m of ward boundaries. Controls exclude non-switching blocks in origin wards with 2014--2015 aldermanic turnover.}
    \label{fig:treatment_control_maps}
\end{figure}

\subsection{Estimating Equation}

The outcome is the count of high-discretion permits issued on each census block in each year. Because permits are counts that aggregate naturally to a balanced block-year panel, with every block observed in every year (including zeros), I can include block fixed effects and identify entirely off within-block variation over time. I estimate the event study using Poisson pseudo-maximum likelihood (PPML), with the log conditional mean given by:
%
\begin{equation}
\label{eq:event_study_permits}
\log E[Y_{bt}] = \alpha_b + \lambda_{gt} + \sum_{\tau \neq -1} \beta_\tau \left( \Delta\text{Stringency}_b \times \mathbf{1}[t = \tau] \right) + \theta_t P_b^{pre} + \psi_t \mathbf{1}[P_b^{pre}=0]
\end{equation}
%
where $Y_{bt}$ is the count of high-discretion permits issued on block $b$ in year $t$, $\alpha_b$ is a block fixed effect, and $\lambda_{gt}$ is a border-pair-by-year fixed effect. In the compiled paper's 2015-cohort design, the treatment variable $\Delta\text{Stringency}_b$ is the destination alderman's baseline strictness score minus the origin alderman's baseline strictness score, where origin and destination aldermen are assigned using the 2014 pre-redistricting ward-to-alderman mapping. The terms $P_b^{pre}$ and $\mathbf{1}[P_b^{pre}=0]$ are the block's raw count of issued high-discretion permits in the pre-redistricting years and an indicator for having no such permits; both are interacted with calendar year to allow blocks with different baseline development intensity to follow different time profiles.

The block fixed effects absorb all time-invariant characteristics of each block, so identification comes from within-block changes over time: the same physical block before and after its alderman changes. The border-pair-by-year fixed effects absorb local shocks that affect all blocks near a given ward boundary in a given year, so the comparison is between blocks near the same boundary that experienced different changes in aldermanic representation. The sample uses the 2015 redistricting cohort, a 250m bandwidth around ward boundaries, and uniform weights. I use a wider window here than in Section~\ref{sec:empirical_results_density} because this design does not identify from a discontinuity exactly at the boundary. It compares treated and control blocks within the same boundary neighborhoods over time, so the extra width preserves local support and power without changing the fixed-effects comparison. Blocks with zero permits in some years remain in the panel and are informative under PPML. Blocks with zero permits in every sample year contribute no within-block variation once block fixed effects are included, so the estimator drops them mechanically. Standard errors are clustered at the census block level.

This within-block design is not feasible for the home sales analysis in Section~\ref{sec:empirical_results_home_sales}. Home sales are individual transactions with property-level hedonic characteristics that vary across properties within the same block, and many blocks have zero or one sale in a given year, which makes block-year aggregation uninformative and block fixed effects impractical at the transaction level. The permit count outcome, by contrast, aggregates cleanly to a balanced panel in which block fixed effects absorb permanent differences in development activity across locations.

\subsection{Results}

Figure~\ref{fig:permit_event_study} presents the main result. The pre-period coefficients are statistically indistinguishable from zero, and the joint pre-trend test does not reject. At redistricting ($t = 0$), high-discretion permit issuance drops sharply, and the effect remains negative through most of the post-period. The pattern is consistent with an immediate disruption to the permitting pipeline when the regulatory gatekeeper changes.

\begin{figure}[htbp]
    \centering
    \includegraphics[width=0.7\textwidth]{../tasks/run_event_study_permit/output/event_study_yearly_cohort_2015_high_discretion_issue_ppml_continuous_uniform_250m_within_block_full_clust_block_ctrl_pre_high_level_geo_wardpair.pdf}
    \caption{Permit Event Study: Effect of Alderman Stringency on High-Discretion Permits}
    \figurenotes{Block-year PPML event study. Outcome is the count of high-discretion permits issued per block per year. 2015 redistricting cohort; sample restricted to blocks within 250m of ward boundaries with uniform weights. Continuous treatment: effect of a 1 SD increase in alderman stringency. Block and border-pair $\times$ year fixed effects; pre-period high-discretion permit count and no-pre-period high-discretion permit indicator interacted with calendar year. Year $-1$ omitted as reference. 95\% confidence intervals based on standard errors clustered at the block level. Pooled post-period estimate: $-0.079$ (SE $= 0.041$). Joint pre-trend test $p = 0.792$.}
    \label{fig:permit_event_study}
\end{figure}

Figure~\ref{fig:permit_event_study_split} traces the effect separately for blocks that moved to more stringent versus more lenient aldermen. The plotted coefficients are directionally consistent with the baseline specification: blocks assigned to more stringent aldermen tend to show lower permit issuance, while blocks assigned to more lenient aldermen tend to move in the opposite direction. I treat this split figure as descriptive rather than as a standalone formal test.

\begin{figure}[htbp]
    \centering
    \includegraphics[width=0.7\textwidth]{../tasks/run_event_study_permit/output/event_study_yearly_cohort_2015_high_discretion_issue_ppml_continuous_split_uniform_250m_within_block_full_clust_block_ctrl_pre_high_level_geo_wardpair.pdf}
    \caption{Permit Event Study: Continuous Split by Treatment Direction}
    \figurenotes{Same specification as Figure~\ref{fig:permit_event_study}, estimated separately for blocks with $\Delta\text{Stringency} > 0$ (moved to more stringent) and $\Delta\text{Stringency} < 0$ (moved to more lenient). 95\% confidence intervals based on standard errors clustered at the block level.}
    \label{fig:permit_event_study_split}
\end{figure}

Table~\ref{tab:did_permit_2015_current} collapses the post-period into a single pooled DID coefficient. The pooled PPML estimate is $-0.079$ (SE $= 0.041$), which is about an 8\% decline for a one standard deviation increase in alderman stringency. This is a direct measure of how aldermanic discretion affects the development pipeline: more stringent aldermen do not simply alter the scale of completed buildings, they reduce the flow of projects that require their engagement.

\begin{table}[htbp]
    \caption{Effect of Alderman Stringency on High-Discretion Permits}
    \label{tab:did_permit_2015_current}

    \input{../tasks/run_event_study_permit/output/did_table_permit_2015_high_discretion_issue_ppml_uniform_250m_ctrl_pre_high_level_geo_wardpair.tex}

    \footnotesize
    \textit{Notes:} Block-year PPML regression using the 2015 implementation cohort, restricted to census blocks within 250m of ward boundaries and relative years $-5$ through $5$. Outcome is the number of issued high-discretion permits. Uniform weighting. The pre-period permit controls are the block's raw count of issued high-discretion permits before redistricting and an indicator for no pre-period high-discretion permits, each interacted with calendar year. Standard errors clustered at the block level. * $p<0.10$, ** $p<0.05$, *** $p<0.01$.
\end{table}

As a placebo, I run the same specification on low-discretion permits, which are less exposed to aldermanic intervention than the high-discretion permits that underlie the stringency index. Figure~\ref{fig:permit_event_study_low_disc} shows the corresponding event study. The plot does not show a clear post-treatment pattern. This makes it less likely that the high-discretion result is picking up a general shift in permitting activity unrelated to discretion.

\begin{figure}[htbp]
    \centering
    \includegraphics[width=0.7\textwidth]{../tasks/run_event_study_permit/output/event_study_yearly_cohort_2015_low_discretion_nosigns_issue_ppml_continuous_uniform_250m_within_block_full_clust_block_ctrl_pre_high_level_geo_wardpair.pdf}
    \caption{Placebo: Low-Discretion Permits}
    \figurenotes{Same specification as Figure~\ref{fig:permit_event_study} but using low-discretion permits (excluding signs) as the outcome. The plotted coefficients are small and do not show a clear post-treatment pattern. The absence of a placebo response suggests that the main result is specific to the discretionary part of the permitting process.}
    \label{fig:permit_event_study_low_disc}
\end{figure}

\subsection{Interpretation and Additional Checks}

Because the specification uses block fixed effects, the design is less exposed to time-invariant differences across blocks that could bias the border discontinuity estimates.

The permits studied here are also not simply inputs to the stringency index. The stringency score is constructed from processing \textit{times} on all high-discretion permits across a ward, while the event study outcome is the \textit{count} of high-discretion permits issued on individual blocks. The score measures how an alderman handles permits in general; the event study asks whether the volume of development-related permitting changes when that alderman gains jurisdiction over new territory.
